\chapter[APRESENTAÇÃO DA TRAJETÓRIA DO ALUNO]{APRESENTAÇÃO DA TRAJETÓRIA DO ALUNO}

Apesar de sempre ter sido apaixonada por teconologia e me interessar pela possibilidade de atuação profissional nessa área, gosto de pensar que minha trajetória profissional tem seu ponto de partida datado em 2020. Nesse ano, parti em uma viagem de intercâmbio para Vancouver, Canadá com o principal objetivo de aprimorar meus conhecimentos de língua inglesa. Durante minha passagem pelo país, tive a oportunidade de conhecer diversas universidades, com destaque para a LaSalle College Vancouver que me proporcionou a experiência de participar de uma oficina de Game Design. Tal experiência culminou no desenvolvimento de um sonho que seria cultivado ao longo de meus próximos três anos de ensino médio: atuar profissionalmente na área de Tecnologia de Informação. 

Desde então, busquei aprimorar meu domínio linguístico, recebendo meu certificado de conclusão do programa W-A-Y American School em 2021, comprovando minha fluência em inglês e, desde 2022, venho me dedicando continuamente ao aprendizado da língua chinesa. Além disso, participei de eventos de Open Campus, escolhendo, enfim, a \ac{pucrs} como minha instituição de ensino para os próximos anos. 

No terceiro semestre, tive a oportunidade de aplicar os conhecimentos adquiridos até então, participando do projeto "Turismo Rural em São Francisco de Paula" por meio da \ac{ages}, o que agregou valiosas experiências nos âmbitos acadêmico e profissional. No final desse semestre, fui selecionada em um processo seletivo de estágio na área de pesquisa e desenvolvimento de software para sistemas aviônicos e defesa militar, onde aprendi e desenvolvi habilidades com HTML, CSS, JavaScript, TypeScript e React.js.

No quarto semestre, dei continuidade à minha trajetória na \ac{ages}, contribuindo para o desenvolvimento do projeto "meditAMAmente".

Atualmente, no quinto semestre, minha trajetória na \ac{ages} segue com o projeto SalesPilot, no qual atuo como \ac{ages} III, assumindo a responsabilidade pela arquitetura do frontend da plataforma. Essa experiência, detalhada ao longo deste documento, tem sido especialmente significativa para meu desenvolvimento profissional, por exigir não apenas a aplicação de conhecimentos técnicos consolidados ao longo do curso, mas também o exercício de um papel de liderança técnica e de comunicação dentro de um time multidisciplinar.